\documentclass[a4paper]{article}
% Español (México) con XeLaTeX: fontspec para fuentes Unicode, polyglossia
% para la división silábica y los nombres del documento en la variante mexicana.
\usepackage{fontspec}
\usepackage{polyglossia}
\setdefaultlanguage[variant=mexican]{spanish}
% Los nombres chinos van como «pinyin (original)»: xeCJK da a los caracteres
% Han su propia fuente; sin ella XeLaTeX los omite con un aviso.
% FandolSong no cubre algunos caracteres de nombres (U+73FA); WenQuanYi Zen Hei sí.
\usepackage[AutoFallBack=true]{xeCJK}
\setCJKmainfont{FandolSong-Regular.otf}
\setCJKfallbackfamilyfont{\CJKrmdefault}{WenQuanYi Zen Hei}
% polyglossia no traduce los nombres de listings.
\AtBeginDocument{\providecommand{\lstlistingname}{}\renewcommand{\lstlistingname}{Listado}%
  \providecommand{\lstlistlistingname}{}\renewcommand{\lstlistlistingname}{Índice de listados}}
\usepackage{amsmath,amssymb}
\usepackage{graphicx}
\usepackage[margin=2.35cm]{geometry}
\usepackage[most]{tcolorbox}
\usepackage{etoolbox}
\usepackage{listings}
\usepackage{xcolor}
\usepackage{hyperref}
\usepackage{booktabs,longtable,tabularx,array}
\usepackage{subcaption}
\usepackage{float}
\usepackage{enumitem}
\usepackage{microtype}
\usepackage{fancyhdr}
\usepackage{needspace}

\definecolor{kbBack}{HTML}{F2F6FA}
\definecolor{kbFrame}{HTML}{9DB4CC}
\definecolor{kbTitle}{HTML}{52708F}
\definecolor{ibBack}{HTML}{FBF5E7}
\definecolor{ibFrame}{HTML}{C9A961}
\definecolor{ibTitle}{HTML}{7D5F1E}
\definecolor{wbBack}{HTML}{FCF2F0}
\definecolor{wbFrame}{HTML}{D6A096}
\definecolor{wbTitle}{HTML}{9E4F43}
\definecolor{dbBack}{HTML}{F4F7F3}
\definecolor{dbFrame}{HTML}{AEC2AC}
\definecolor{dbTitle}{HTML}{5F7F64}

\tcbset{softbox/.style={
  enhanced,breakable,boxrule=0.8pt,arc=2pt,left=3mm,right=3mm,
  top=2mm,bottom=2mm,coltitle=white,fonttitle=\bfseries,
  toptitle=0.6mm,bottomtitle=0.6mm
}}
\newtcolorbox{knowledgebox}[1]{softbox,colback=kbBack,colframe=kbFrame,colbacktitle=kbTitle,title={#1}}
\newtcolorbox{importantbox}[1]{softbox,colback=ibBack,colframe=ibFrame,colbacktitle=ibTitle,title={#1}}
\newtcolorbox{warningbox}[1]{softbox,colback=wbBack,colframe=wbFrame,colbacktitle=wbTitle,title={#1}}
\newtcolorbox{dialoguebox}[1]{softbox,colback=dbBack,colframe=dbFrame,colbacktitle=dbTitle,title={#1}}

\lstset{
  language=Python,basicstyle=\ttfamily\small,
  keywordstyle=\color{kbTitle}\bfseries,stringstyle=\color{wbTitle},
  commentstyle=\color{dbTitle},breaklines=true,frame=single,numbers=left,
  numberstyle=\tiny\color{gray},captionpos=b,columns=fullflexible,
  keepspaces=true,showstringspaces=false,extendedchars=true
}
\BeforeBeginEnvironment{lstlisting}{\Needspace{12\baselineskip}}

\hypersetup{colorlinks=true,linkcolor=kbTitle,urlcolor=kbTitle,citecolor=wbTitle}
\setlist{itemsep=0.25em,topsep=0.4em}
\setlength{\parindent}{2em}
\setlength{\parskip}{0.25em}
\newcolumntype{L}[1]{>{\raggedright\arraybackslash}p{#1}}

\providecommand{\notetitle}{Notas del curso: [tema del curso]}
\providecommand{\noteauthors}{Elaboradas a partir de los materiales públicos de Stanford CS25}
\providecommand{\notedate}{Versión reescrita en 2026}
\providecommand{\videochannel}{Stanford Online}
\providecommand{\videopublishdate}{2022}
\providecommand{\videoduration}{[duración del video]}
\providecommand{\videourl}{}
\providecommand{\videocoverpath}{cover.jpg}
\providecommand{\coursesite}{https://web.stanford.edu/class/cs25/}
\providecommand{\repourl}{https://github.com/hqhq1025/ai-course-notes}
\providecommand{\skillversion}{wdkns/wdkns-skills@39f1a04}

\newcommand{\figsource}[1]{\par\vspace{-0.3em}{\footnotesize\color{gray}\emph{Fuente: #1}}\par}
\newcommand{\teachervoice}[1]{\begin{knowledgebox}{Nota de clase}#1\end{knowledgebox}}
\newcommand{\term}[1]{\textbf{#1}}

\pagestyle{fancy}
\fancyhf{}
\fancyfoot[C]{\thepage}
\fancyfoot[R]{\footnotesize\color{gray}\href{\repourl}{AI Course Notes}}
\renewcommand{\headrulewidth}{0pt}
\renewcommand{\footrulewidth}{0pt}

\newcommand{\makecscover}{
\begin{titlepage}
\centering
\vspace*{0.7cm}
{\Large Stanford CS25 · Transformers United\par}
\vspace{0.8cm}
{\huge\bfseries \notetitle\par}
\vspace{0.6cm}
{\large \noteauthors\par}
\vspace{0.25cm}
{\large \notedate\par}
\vspace{0.9cm}
\ifdefempty{\videocoverpath}{}{
  \includegraphics[width=0.86\textwidth,height=0.43\textheight,keepaspectratio]{\videocoverpath}\par
}
\vfill
\begin{tcolorbox}[width=0.92\textwidth,colback=black!2!white,colframe=black!45,boxrule=0.7pt,arc=2pt]
\textbf{Autor o canal del video}: \videochannel\par
\textbf{Fecha de publicación}: \videopublishdate\par
\textbf{Duración del video}: \videoduration\par
\textbf{Sitio del curso}: \href{\coursesite}{\nolinkurl{\coursesite}}\par
\ifdefempty{\videourl}{}{\textbf{Enlace del video}: \href{\videourl}{\nolinkurl{\videourl}}\par}
\textbf{Normas de elaboración}: \skillversion\ + las normas source-first / teacher-voice / visual-QA de este repositorio
\end{tcolorbox}
\end{titlepage}
\tableofcontents
\newpage
}

\usepackage{multicol}

\renewcommand{\notetitle}{Lecture 33: Comprender la fuerza dominante de la IA a partir de la historia del Transformer}
\renewcommand{\noteauthors}{Elaborado a partir de la conferencia independiente oficial de Hyung Won Chung, el deck de 67 páginas y los subtítulos hechos por personas}
\renewcommand{\notedate}{CS25 V4 · versión reescrita source-first de 2026}
\renewcommand{\videochannel}{Stanford Online}
\renewcommand{\videopublishdate}{Clase: 2024-04-11; publicación: 2024-06-11}
\renewcommand{\videoduration}{36:30}
\renewcommand{\videourl}{https://www.youtube.com/watch?v=orDKvo8h71o}
\renewcommand{\videocoverpath}{cover.jpg}

\newcommand{\lecturefigure}[3]{%
\begin{figure}[H]
  \centering
  \includegraphics[width=0.97\textwidth,height=0.49\textheight,keepaspectratio]{slides-images/#1}
  \caption{#2}
  \figsource{#3}
\end{figure}
}

\let\standardtableofcontents\tableofcontents
\renewcommand{\tableofcontents}{%
  \begingroup
  \small
  \setlength{\parskip}{0pt}
  \standardtableofcontents
  \endgroup
}

\begin{document}
\makecscover

\section{Estudiar el cambio mismo: identificar primero la fuerza dominante}

Cuando la IA avanza demasiado rápido, seguir las noticias una por una acumula hechos, pero no necesariamente permite formarse un juicio sobre la dirección. Hyung Won Chung elige otro camino: trata la historia de la arquitectura Transformer como un conjunto de experimentos naturales y se pregunta qué diseños eran razonables con las tareas y la capacidad de cómputo de antes, y cuáles perdieron su ventaja al cambiar la escala, los datos y las formas de interacción. El objetivo de la clase no es demostrar que decoder-only siempre es correcto, sino entrenar un método de investigación transferible.

\subsection{Por qué el análisis histórico no es nostalgia}

La portada coloca ``history'' y ``future'' en la misma oración. Al leer las dos diapositivas siguientes conviene captar primero la metodología: el objeto de estudio no es el nombre de algún modelo, sino la trayectoria del cambio. Si se identifican las variables direccionales que gobiernan la trayectoria, las actualizaciones dispersas pueden condensarse en unas pocas hipótesis verificables. Cada comparación de arquitecturas que sigue debe volver a esta pregunta: ¿registra a un ganador estático o explica por qué la ventaja relativa se desplaza con el regime?

\lecturefigure{hyung-slide-001.jpg}{Tema de la conferencia: moldear el futuro de la IA a partir de la historia del Transformer}{Hyung Won Chung official deck, p.~1}

\lecturefigure{hyung-slide-002.jpg}{No basta con seguir los resultados más recientes: hay que estudiar el cambio mismo}{Hyung Won Chung official deck, p.~2}

\teachervoice{00:02:05--00:02:58, Hyung Won enfatiza: cuando la velocidad del cambio supera la capacidad humana de seguirlo, hay que mirar los sistemas anteriores y trazar el camino de «cómo llegamos hasta aquí». Lo central de la clase no es memorizar la tabla de clasificación de 2024, sino abstraer la dirección a partir del camino.}

\begin{warningbox}{Límites de la evidencia de la edición independiente}
Este video es una edición independiente de 36:30 de la lista de reproducción oficial; su contenido coincide con la segunda mitad de la intervención de Hyung Won en la grabación conjunta de la Lecture 27, pero no incluye la Q\&A conjunta posterior. Estas notas conservan solo la prepared talk que aparece efectivamente en el video independiente y no incorporan la discusión de la Q\&A sobre MLE, RLHF o los límites energéticos.
\end{warningbox}

\subsection{Identify, understand, roll forward}

La tercera diapositiva divide la predicción en tres pasos: primero identificar la dominant driving force, luego entender cómo modifica las ventajas y desventajas relativas entre métodos y, por último, extrapolar la tendencia con cautela. La fuerza dominante no es la única fuerza, sino la variable que explica la mayor parte de la direccionalidad con la precisión que se requiere en ese momento; si el regime cambia, hay que volver a modelar. Antes de leer la figura también conviene distinguir entre «describir el pasado» y «predecir el futuro»: lo primero puede acotarse con datos históricos; lo segundo debe declarar a la vez sus condiciones de falla y sus explicaciones alternativas.

\lecturefigure{hyung-slide-003.jpg}{Método de tres pasos para estudiar el cambio: identificar, entender y extrapolar la fuerza dominante}{Hyung Won Chung official deck, p.~3}

\begin{importantbox}{Definición operativa de dominant driving force}
La fuerza impulsora dominante es el factor que basta para explicar los principales cambios de dirección en una escala de tiempo, un presupuesto y una tolerancia al error dados. Permite al investigador ignorar temporalmente los efectos menores a cambio de poder analizar el problema; pero lo que se ignora debe declararse, y cuando la predicción falla, lo primero que hay que revisar son las retroalimentaciones y restricciones omitidas.
\end{importantbox}

\teachervoice{00:03:03--00:04:07, el ponente resume el proceso de investigación como identify, understand, roll forward, y señala que una predicción no necesita una tasa de acierto alta para ser valiosa. Lo importante es que las hipótesis puedan verificarse, no fingir que se conoce el futuro con precisión.}

\subsection{El experimento de la pluma que cae: cuándo puede ignorarse la resistencia del aire}

Cuando se suelta una pluma desde el reposo, en la realidad actúan a la vez la gravedad, la resistencia del aire, la rotación y la deformación del material; aun así, en la clase solo se modela la gravedad, porque en una distancia corta y con una precisión ordinaria es suficientemente dominante. Si se toma como positiva la dirección hacia abajo, puede escribirse
\[
y(t)=y_0+v_0t+\frac{1}{2}gt^2.
\]
donde $y(t)$ es la posición en el instante $t$, $y_0$ es la posición inicial, $v_0$ es la velocidad inicial y $g$ es la aceleración de la gravedad. La condición para que la fórmula sea válida no es que «la resistencia no exista», sino que «la contribución de la resistencia al error sea aceptable para el problema en cuestión».

\lecturefigure{hyung-slide-004.jpg}{El experimento de la pluma que cae: la gravedad ilustra cómo ignorar las fuerzas secundarias}{Hyung Won Chung official deck, p.~4}

\lecturefigure{hyung-slide-006.jpg}{Cuantas más fuerzas dominantes haya y más complejas sean sus interacciones, más difícil es predecir la trayectoria}{Hyung Won Chung official deck, p.~6}

\begin{knowledgebox}{Lectura de la figura: el eje horizontal no es el número de campos observados}
El eje horizontal representa el \textbf{número de fuerzas dominantes en competencia}. Un sistema puede tener muchas características y, aun así, su dirección puede estar determinada por unos pocos mecanismos; a la inversa, con pocas variables pero retroalimentaciones fuertemente acopladas, también puede ser difícil extrapolar. Lo que la figura respalda es un orden de modelado: primero encontrar el factor de mayor direccionalidad y después decidir cuándo incorporar las interacciones.
\end{knowledgebox}

\teachervoice{00:04:40--00:05:57, el ponente menciona explícitamente la resistencia del aire y las interacciones aerodinámicas complejas, pero decide ignorarlas. La nota de clase no es que «la realidad tiene una sola variable», sino «primero se define la precisión requerida y luego se decide qué variables pueden aproximarse».}

\subsection{Resumen de la sección}

Esta sección establece la disciplina de evidencia de toda la clase: la historia no es una lista de modelos, sino un experimento natural; la fuerza dominante no es la única causa, sino una aproximación con un margen de error declarable; la extrapolación no es una garantía, sino la forma de generar el siguiente conjunto de experimentos. A continuación, este método se aplica a compute, inductive bias y la arquitectura Transformer.
\section{Compute, Bitter Lesson y el ciclo de vida del sesgo inductivo}

La introducción ya presentó el método de «identificar la fuerza dominante»; ahora se aplica a la historia de la investigación en IA. El juicio central de Hyung Won es que la disminución sostenida del costo de cómputo y el associated scaling constituyen una fuerza direccional importante; por ello, el diseño de una estructura no debe evaluarse solo por si hoy es más preciso, sino también por si limitará la posibilidad de seguir aumentando la escala mañana. Lo importante aquí no es afirmar una causalidad única, sino volver a comparar las decisiones de arquitectura dentro de un compute regime.

\subsection{Por qué tomar el compute per dollar como fuerza dominante}

Esta sección comienza por la evidencia. El eje vertical de la figura no es el número de transistores, sino la capacidad de cómputo que puede adquirirse con un monto fijo de dinero; el eje horizontal abarca más de un siglo. A partir de esta tendencia de largo plazo, Hyung Won deriva una estrategia de investigación: no competir con una tendencia externa poderosa y persistente, sino diseñar métodos capaces de aprovecharla. Al mismo tiempo, es necesario incluir el hardware, los algoritmos, la precisión numérica y la eficiencia del software dentro del «cómputo efectivo»; de lo contrario, se confundiría el indicador de un solo componente con la trayectoria completa de la investigación.

\lecturefigure{hyung-slide-007.jpg}{La capacidad de cómputo obtenible con un costo fijo crece a largo plazo, aproximadamente diez veces cada cinco años}{Hyung Won Chung official deck, p.~7}

Si la capacidad de cómputo aumenta aproximadamente diez veces cada cinco años, puede escribirse
\[
C(t)=C_0\,10^{t/5},
\]
donde $C(t)$ es la cantidad de cómputo que puede comprarse con un presupuesto fijo $t$ años después del punto de partida, y $C_0$ es la cantidad de cómputo en el punto de partida. Esta expresión es solo una aproximación de la figura de la clase y no garantiza que en el futuro se mantenga el mismo exponente; la energía, la manufactura, la interconexión, el capital y los algoritmos pueden modificar la tasa de crecimiento efectiva.

\begin{knowledgebox}{Lectura de la figura: el compute per dollar no es un indicador de hardware aislado}
La figura de la clase agrega el efecto de largo plazo del hardware, los algoritmos, la precisión numérica, la pila de software y la integración de sistemas. La pregunta pertinente es «cuánto cómputo de entrenamiento efectivo puede realizarse con un costo fijo», no solo seguir la cantidad de transistores. La figura muestra una tendencia histórica; no demuestra que las restricciones de energía, manufactura, interconexión y capital no vayan a cambiar la pendiente en el futuro.
\end{knowledgebox}

\teachervoice{La formulación de 00:08:00--00:09:03 es «no compitan con esta tendencia; procuren aprovecharla». Se trata de una heurística para elegir líneas de investigación, no de una garantía de crecimiento exponencial ilimitado; por eso, estas notas separan la observación histórica, la estrategia de investigación y la predicción del futuro.}

\subsection{Bitter Lesson: imponerle menos a la máquina la forma de pensar que creemos entender}

La sección anterior mostró que la escala es una tendencia fuerte; esta explica por qué el exceso de estructura diseñada a mano entra en conflicto con ella. Los investigadores suelen codificar en el modelo su propia explicación del pensamiento humano; con poco cómputo, estas estructuras pueden parecer atajos, pero restringen la exploración de otras representaciones por parte del modelo. La versión de la Bitter Lesson que se presenta en clase es: el progreso de largo plazo proviene una y otra vez de métodos más generales, con supuestos más débiles, sumados a más datos y más cómputo.

\lecturefigure{hyung-slide-008.jpg}{«Enseñar a la máquina a pensar como creemos que pensamos» puede imponer una estructura equivocada}{Hyung Won Chung official deck, p.~8}

\lecturefigure{hyung-slide-009.jpg}{Bitter Lesson: métodos más generales, supuestos de modelado más débiles y mayor escala}{Hyung Won Chung official deck, p.~9}

\begin{importantbox}{Qué es el inductive bias}
El inductive bias (sesgo inductivo) es el supuesto estructural por el cual un modelo, con datos limitados, prefiere cierta clase de funciones, representaciones o formas de interacción. La localidad de la convolución, la separación de parámetros entre encoder y decoder, la attention bidireccional y ciertas invariancias específicas son todos sesgos. Un sesgo puede mejorar la eficiencia muestral, pero también excluye parte de las soluciones; su valor depende de los datos, el cómputo, la tarea y las restricciones de despliegue vigentes.
\end{importantbox}

\teachervoice{En 00:09:17--00:11:19, Hyung Won afirma que los investigadores con frecuencia intentan modelar «cómo creemos que pensamos», sin entender realmente los mecanismos cognitivos de bajo nivel. La postura de la clase es deliberadamente tajante; estas notas la conservan como una advertencia de diseño, no como prueba de que toda estructura previa sea errónea.}

\subsection{Menos estructura es más scalable, pero no necesariamente mejor hoy}

Si se compara la cantidad de estructura de forma transversal, los métodos con más estructura suelen partir de un nivel más alto en el intervalo de bajo cómputo, pero pueden saturarse antes; los métodos con menos estructura son difíciles de entrenar al principio, pero le dan al modelo más grados de libertad y, conforme crece la escala, pueden superar a los primeros. Aquí, ``scalable'' significa que la curva de crecimiento puede seguir mejorando, no solo que el método pueda ejecutarse en más GPU.

\lecturefigure{hyung-slide-010.jpg}{Más estructura mejora el desempeño con poco cómputo, pero puede reducir la escalabilidad a largo plazo}{Hyung Won Chung official deck, p.~10}

\begin{knowledgebox}{Lectura de la figura: comparar el punto de cruce, no etiquetar la estructura como buena o mala}
A la izquierda del presupuesto actual, el método azul puede ser claramente mejor; el método rojo solo lo supera con un compute mayor. La decisión de investigación depende de dónde se ubica el punto de despliegue, de si puede asumirse el costo de exploración y de si en el futuro realmente se alcanzará el punto de cruce. La figura respalda una regime-dependent choice, no la idea de «elegir siempre la menor estructura».
\end{knowledgebox}

\lecturefigure{hyung-slide-013.jpg}{En el compute regime actual conviene elegir un sesgo que funcione, sin olvidar retirarlo en el futuro}{Hyung Won Chung official deck, p.~13}

\lecturefigure{hyung-slide-014.jpg}{Dados el cómputo, los datos, los algoritmos y la arquitectura, existe un inductive bias óptimo de manera temporal}{Hyung Won Chung official deck, p.~14}

Si se denota la intensidad de la estructura como $b$ y el estado actual de los recursos como $r$, el desempeño puede escribirse como
\[
b^{\star}(r)=\arg\max_b\; P(b;r).
\]
Aquí, $P(b;r)$ es el desempeño del sesgo $b$ en el estado de recursos $r$, y $b^{\star}(r)$ es la estructura óptima actual. A medida que cambian el compute, los data y los algorithm que contiene $r$, $b^{\star}$ también se desplaza; por lo tanto, una decisión de arquitectura que alguna vez fue correcta no debe congelarse de manera permanente.

\teachervoice{La salvedad principal de 00:12:04--00:13:27 es que no se puede esperar desde el inicio al método más general, porque podría no funcionar en absoluto con el presupuesto actual. La acción correcta es «incorporar ahora el sesgo útil y registrar explícitamente cuándo revisarlo y retirarlo en el futuro».}

\subsection{Por qué los incentivos de investigación favorecen agregar estructura y no eliminarla}

El ciclo de vida de la estructura también tiene causas sociales. Agregar un módulo, un término de pérdida o un mecanismo se convierte fácilmente en la contribución de un artículo; demostrar que un módulo complejo puede eliminarse a mayor escala suele parecer poco novedoso. Por eso, la comunidad acumula supuestos diseñados para regímenes anteriores, pero carece de una sustracción sistemática. Esta sección considera también los incentivos de investigación como parte de la evolución de la arquitectura, porque qué experimentos reciben recursos y qué resultados negativos se publican influye directamente en cuánto tiempo se conservan los sesgos antiguos.

\lecturefigure{hyung-slide-015.jpg}{El método que es mejor a largo plazo con frecuencia es peor en el presente}{Hyung Won Chung official deck, p.~15}

\begin{warningbox}{Una ventaja de corto plazo en el leaderboard no equivale a valor de investigación de largo plazo}
Si un método solo lleva la delantera con modelos pequeños, pocos datos o un benchmark específico, hay que seguir preguntando: ¿la ventaja proviene de un supuesto previo más fuerte? ¿Se satura al aumentar la escala? ¿La complejidad de implementación dificulta entrenamientos más grandes? Estas preguntas no pueden responderse con la clasificación puntual actual.
\end{warningbox}

\teachervoice{En 00:13:27--00:14:23, Hyung Won dice sin rodeos que los incentivos de publicación fomentan agregar estructura, pero rara vez premian eliminarla. Añade que el método que es mejor a largo plazo casi inevitablemente se ve peor hoy, porque un sistema de aprendizaje con más libertad es más caótico en el intervalo de pocos recursos.}

\lecturefigure{hyung-slide-016.jpg}{Síntesis de la etapa: la fuerza dominante es un cómputo más barato y el associated scaling}{Hyung Won Chung official deck, p.~16}

Esta diapositiva de síntesis no descarta todas las demás fuerzas, sino que anuncia el sistema de referencia del análisis que sigue: interpretar, una por una, las estructuras adicionales de la arquitectura Transformer como un inductive bias propio de ciertas condiciones históricas, y luego preguntar si la escala actual todavía lo necesita. Así, el diagrama de arquitectura deja de ser una clasificación estática y se convierte en un objeto sobre el que puede razonarse a lo largo de la scale trajectory.

\subsection{Resumen de la sección}

Esta sección estableció el marco de juicio de Hyung Won: la disminución sostenida del costo de cómputo les da a los métodos generales un gran potencial a largo plazo; el sesgo inductivo mejora la eficiencia cuando los recursos son escasos, pero puede limitar el escalamiento posterior. La conclusión correcta no es «cuanta más estructura, mejor» ni «cuanta menos estructura, mejor», sino estimar de manera continua el sesgo óptimo del regime actual y dejar abierto el camino para futuros experimentos de sustracción.
\section{Tres tipos de Transformer: a partir de la base común del sequence model}

El capítulo anterior discutió por qué conviene estudiar la estructura; este capítulo primero la explica con claridad. Encoder-decoder, encoder-only y decoder-only no son tres invenciones sin relación entre sí: las tres convierten una entrada discreta en una secuencia de vectores y después usan attention para modelar la interacción entre los elementos. Las diferencias provienen principalmente de si los parámetros se comparten, de cómo la attention mask limita el alcance visible y de cómo se entrena y se usa la salida.

\subsection{La interfaz de tarea de las tres arquitecturas}

Esta sección presenta primero el panorama general y después entra paso a paso en las figuras. Encoder-decoder asigna dos stacks distintos a la entrada y a la salida; encoder-only comprime la entrada en una representación y por lo general requiere un task-specific head; decoder-only concatena la entrada y el objetivo, y genera token por token en un único causal stack. Cada interfaz implica supuestos distintos sobre la tarea.

\lecturefigure{hyung-slide-017.jpg}{Los tres tipos de arquitectura Transformer: encoder-decoder, encoder-only y decoder-only}{Hyung Won Chung official deck, p.~17}

\begin{knowledgebox}{Glosario: qué problema resuelve cada arquitectura}
\begin{tabularx}{\textwidth}{L{0.22\textwidth} X}
\toprule
Arquitectura & Mecanismo central y relación con el curso \\
\midrule
Encoder-decoder & El encoder codifica la entrada de forma bidireccional; el decoder genera el objetivo mediante causal self-attention y cross-attention. Es adecuada para sequence-to-sequence explícito. \\
Encoder-only & Interacción bidireccional sobre toda la secuencia; la salida suele usarse para clasificación, recuperación o token labeling. Para generar una secuencia necesita un mecanismo de generación adicional. \\
Decoder-only & La entrada y la salida previa están en la misma secuencia causal; comparte parámetros y tiene una interfaz unificada, por lo que es adecuada de forma natural para autoregressive generation y para la continuación en varios turnos. \\
\bottomrule
\end{tabularx}
\end{knowledgebox}

\subsection{De los caracteres a la secuencia de vectores, y de ahí a la interaction}

Antes de clasificar las arquitecturas hay que entender el pipeline común. Primero, el tokenizer asigna al texto token IDs; después, la embedding table asigna a cada ID un vector de dimensión $d_{\text{model}}$; por último, el sequence model permite que las distintas posiciones intercambien información según el attention pattern. Lo que el modelo procesa es una secuencia de vectores, no directamente el «significado de las palabras».

\lecturefigure{hyung-slide-018.jpg}{La secuencia original de caracteres es el punto de partida del pipeline del Transformer}{Hyung Won Chung official deck, p.~18}

\lecturefigure{hyung-slide-021.jpg}{Pipeline completo: tokenización, embedding y sequence model}{Hyung Won Chung official deck, p.~21}

Si el token ID es $i_t$ y la embedding matrix es $E\in\mathbb{R}^{V\times d}$, entonces
\[
h_t^{(0)}=E[i_t]+p_t.
\]
Aquí $V$ es el tamaño del vocabulario, $d$ es la hidden dimension, $E[i_t]$ es el vector del token número $i_t$, $p_t$ es la representación de posición y $h_t^{(0)}$ es el estado que entra a la primera capa. La tokenización determina la longitud de la secuencia y los límites de la segmentación; el embedding determina la representación continua inicial. Ambos influyen en el costo de los cálculos posteriores.

\teachervoice{Entre 00:16:11 y 00:17:26, Hyung Won define primero el Transformer como un sequence model: los elementos de entrada pueden ser palabras, imágenes u otros objetos, y lo central es modelar la interaction entre ellos. Este orden evita separar la fórmula de attention de la interfaz de datos.}

\subsection{Encoder-decoder: tres attention roles}

Ahora pasamos al Transformer original. Al leer la figura, conviene separar primero los tres flujos de información: la self-attention bidireccional dentro del encoder, la self-attention causal dentro del decoder y la cross-attention con la que el decoder consulta la representación del encoder. Cada arista corresponde a un permiso del tipo «qué posición puede leer qué posición».

\lecturefigure{hyung-slide-022.jpg}{Arquitectura encoder-decoder completa y sus tres attention roles}{Hyung Won Chung official deck, p.~22}

\begin{importantbox}{Self-attention y cross-attention}
Dados query $Q$, key $K$ y value $V$, la scaled dot-product attention es
\[
\operatorname{Attn}(Q,K,V)=\operatorname{softmax}\!\left(\frac{QK^{\top}}{\sqrt{d_k}}+M\right)V.
\]
Aquí $d_k$ es la dimensión de key y $M$ es la mask de visibilidad. En la self-attention, $Q,K,V$ provienen de la misma secuencia; en la cross-attention, $Q$ proviene del decoder y $K,V$ provienen de la salida del encoder.
\end{importantbox}

\lecturefigure{hyung-slide-023.jpg}{El encoder usa bidirectional self-attention para modelar la entrada completa}{Hyung Won Chung official deck, p.~23}

La bidirectional attention (atención bidireccional) permite que cualquier posición de la entrada lea cualquier otra posición. Si la longitud de la secuencia es $n$, la mask de visibilidad puede escribirse como $M_{ij}=0$; no hay enmascaramiento del futuro. Es adecuada para una entrada que se proporciona de una sola vez, porque todos los tokens ya se conocen al codificar, pero también implica que un cambio al final de la entrada afecta la representación de todas las posiciones anteriores.

\lecturefigure{hyung-slide-025.jpg}{El decoder usa causal self-attention y solo lee la posición actual y el historial a su izquierda}{Hyung Won Chung official deck, p.~25}

La causal attention (atención causal) usa una mask triangular superior para impedir la lectura del futuro:
\[
M_{ij}=\begin{cases}
0,&j\le i,\\
-\infty,&j>i.
\end{cases}
\]
Aquí $i$ es la posición actual de query y $j$ es la posición de key que se lee. Después de la softmax, el peso de las posiciones futuras se vuelve cero; por eso, durante el entrenamiento pueden calcularse todas las posiciones en paralelo y, durante la inferencia, se mantiene la condición autoregressive.

\lecturefigure{hyung-slide-027.jpg}{La cross-attention conecta las posiciones del objetivo en el decoder con la representación de entrada del encoder}{Hyung Won Chung official deck, p.~27}

\lecturefigure{hyung-slide-028.jpg}{Construcción final del encoder-decoder: codificación bidireccional de la entrada y generación causal del objetivo}{Hyung Won Chung official deck, p.~28}

\begin{knowledgebox}{Lectura de la figura: cómo fluye una ruta de traducción automática}
La entrada ``That is good'' primero interactúa de forma bidireccional dentro del encoder y forma una representación contextualizada de cada token de origen; al generar ``Das ist gut'', el decoder solo puede leer el prefijo del objetivo ya generado y, al mismo tiempo, accede a la última capa del encoder mediante cross-attention. Lo que la figura respalda es la estructura de permisos de información; no indica que los pesos de attention equivalgan por naturaleza a una explicación de alineamiento.
\end{knowledgebox}

\subsection{Diferencias clave entre encoder-only y decoder-only}

La sección anterior usó dos stacks; esta sección examina los dos extremos. Encoder-only produce una representación bidireccional global y, para clasificar, suele leer un token especial y agregar una task-specific layer; decoder-only, en cambio, unifica la entrada y la salida en una sola secuencia y usa el mismo conjunto de parámetros y causal self-attention tanto para comprender el contexto como para generar.

\lecturefigure{hyung-slide-032.jpg}{Forma completa de encoder-only: representación bidireccional, token especial y cabeza de tarea}{Hyung Won Chung official deck, p.~32}

\begin{warningbox}{Encoder-only no significa «no puede producir ninguna secuencia»}
Esta diapositiva se refiere a la interfaz predeterminada: un solo encoder stack no incluye por sí mismo una autoregressive generation path. Es posible agregarle un decoder, un span head, el llenado iterativo de huecos u otro mecanismo de generación; el costo es que el sistema deja de ser la estructura encoder-only mínima de la figura. Lo que el curso compara son los supuestos integrados, no la capacidad absoluta en sentido lógico.
\end{warningbox}

\lecturefigure{hyung-slide-037.jpg}{Forma completa de decoder-only: entrada y objetivo concatenados, parámetros compartidos y causal attention}{Hyung Won Chung official deck, p.~37}

La interfaz unificada de decoder-only puede escribirse como
\[
p(y\mid x)=\prod_{t=1}^{|y|}p(y_t\mid x,y_{<t}).
\]
Aquí $x$ es la secuencia de tokens de entrada, $y$ es la secuencia objetivo y $y_{<t}$ es el prefijo del objetivo ya generado. La entrada y el objetivo usan el mismo stack, pero la loss normalmente puede calcularse solo en las posiciones del objetivo; «compartir parámetros» no significa que el modelo no pueda distinguir los roles, porque la posición, los tokens separadores y el propio contexto siguen proporcionando condiciones.

\subsection{Resumen de la sección}

Los tres tipos de Transformer comparten la base de tokenización, embedding, attention y MLP; las diferencias principales están en la visibilidad de la información, la agrupación de parámetros y la interfaz de tarea. Una vez entendidas estas diferencias, el siguiente capítulo ya no trata encoder-decoder y decoder-only como categorías inconmensurables, sino que transforma de manera continua la primera en la segunda mediante cuatro operaciones, lo que deja al descubierto el verdadero sesgo inductivo.
\section{Transformación en cuatro pasos: reescribir de forma continua un Encoder-decoder como Decoder-only}

La sección anterior presentó los diagramas estáticos de las arquitecturas; esta sección desarrolla la derivación central de Hyung Won. El punto de partida es un encoder-decoder estándar y el punto de llegada es un decoder-only; cada paso elimina una sola diferencia y actualiza la tabla comparativa. El valor didáctico de la derivación no está en recomendar una implementación concreta, sino en reformular «las arquitecturas son distintas» como cuatro hipótesis que pueden someterse a experimentos por separado.

\subsection{Punto de partida: cuatro tipos de diferencias}

Esta subsección fija primero las dimensiones de comparación: la cross-attention adicional, si los parámetros del encoder y del decoder están separados, en qué capa se conecta target-to-input y si la self-attention sobre la entrada es bidireccional. Cada paso posterior cambia solo uno de estos elementos, para no mezclar varios cambios en una misma ablation.

\lecturefigure{hyung-slide-040.jpg}{Punto de partida de la transformación: comparación lado a lado del encoder-decoder estándar y el decoder-only}{Hyung Won Chung official deck, p.~40}

\begin{importantbox}{Cuatro hipótesis verificables}
\begin{enumerate}
  \item ¿La entrada y el objetivo necesitan una attention role independiente?
  \item ¿La entrada y el objetivo necesitan parámetros independientes?
  \item ¿Cada capa del decoder solo puede leer la capa final del encoder?
  \item ¿La entrada debe usar forzosamente bidirectional attention?
\end{enumerate}
Solo al separarlas se puede saber de qué elemento estructural proviene la diferencia de desempeño.
\end{importantbox}

\subsection{Paso uno: compartir los parámetros de cross-attention y self-attention}

Tanto la cross-attention como la self-attention usan query, key, value y output projection; la diferencia principal es de dónde provienen los tensores de entrada. Por eso el primer paso no elimina conexiones, sino que hace que ambas roles compartan los parámetros de proyección, de modo que un solo attention module se encargue a la vez de la interacción dentro de la secuencia y entre secuencias.

\lecturefigure{hyung-slide-041.jpg}{Paso uno: compartir los parámetros de cross-attention y self-attention}{Hyung Won Chung official deck, p.~41}

Si originalmente se tienen $W_Q^{\mathrm{self}},W_K^{\mathrm{self}},W_V^{\mathrm{self}}$ y $W_Q^{\mathrm{cross}},W_K^{\mathrm{cross}},W_V^{\mathrm{cross}}$, compartir parámetros equivale a imponer la restricción
\[
W_{\bullet}^{\mathrm{self}}=W_{\bullet}^{\mathrm{cross}},
\quad \bullet\in\{Q,K,V,O\}.
\]
Aquí $Q,K,V,O$ representan, respectivamente, query, key, value y output projection. El origen de las conexiones puede seguir siendo distinto, pero el modelo ya no puede conservar mapeos lineales completamente independientes para las dos roles.

\lecturefigure{hyung-slide-042.jpg}{Comparación tras el paso uno: la cross-attention independiente pasa a ser una self-attention que cumple ambas roles}{Hyung Won Chung official deck, p.~42}

\begin{knowledgebox}{Lectura de la figura: compartir parámetros no equivale a compartir entradas}
Lo que cambia en la tabla es la module identity, no la visibilidad. El decoder todavía puede aplicar masks distintas a su propio historial y a las representaciones del encoder; lo único es que las matrices de proyección son las mismas. Si en un experimento se cambian a la vez la mask y los parámetros, es difícil determinar si la ganancia proviene de compartir representaciones o del patrón de conexión.
\end{knowledgebox}

\subsection{Paso dos: compartir los parámetros del encoder y del decoder}

El primer paso unificó las attention roles; el segundo unifica las dos stacks. Si tanto la entrada como el objetivo son simplemente token sequences, usar parámetros completamente independientes equivale a suponer que los dos tipos de secuencia requieren espacios de representación distintos. Compartir parámetros, en cambio, exige que un mismo conjunto de layers procese ambas y deja que el contexto y la mask expresen la diferencia. Por eso esta subsección no solo compara el número de parámetros, sino también dos decisiones de modelado: si la diferencia de roles debe escribirse en los pesos o expresarse dinámicamente según la condición de la secuencia.

\lecturefigure{hyung-slide-044.jpg}{Paso dos: compartir los parámetros de las layers del encoder y del decoder}{Hyung Won Chung official deck, p.~44}

\lecturefigure{hyung-slide-045.jpg}{Comparación tras el paso dos: separate parameters pasa a ser shared parameters}{Hyung Won Chung official deck, p.~45}

\begin{warningbox}{Compartir parámetros es un supuesto a priori más débil sobre la tarea y también puede causar interferencia}
Compartir reduce parámetros y favorece la reutilización de conocimiento entre idiomas y entre roles, pero cuando las distribuciones de la entrada y del objetivo difieren mucho, también puede generar competencia por la capacidad. La pregunta correcta no es si «compartir es necesariamente más avanzado», sino si la tarea actual todavía requiere especialización y si esa ganancia de la especialización se mantiene al aumentar la capacidad total.
\end{warningbox}

\teachervoice{La explicación de 00:24:53--00:25:17 es muy directa: un decoder-only tiene una sola stack, así que comparte parámetros de forma natural. Hyung Won usa este contraste para cuestionar si la separación del encoder-decoder sigue siendo necesaria, no para afirmar que compartir no tenga costo.}

\subsection{Paso tres: alinear las representaciones de la misma capa del decoder y del encoder}

En la cross-attention estándar, cada decoder layer suele leer la capa final del encoder. En cambio, la self-attention de misma capa del decoder-only hace que la capa $\ell$ lea las representaciones de los tokens del prefijo en esa misma capa. El tercer paso cambia la conexión a layer-aligned, de modo que la granularidad de las representaciones se acerque más a la de una sola stack.

\lecturefigure{hyung-slide-047.jpg}{Paso tres: decoder layer $\ell$ attends to encoder layer $\ell$}{Hyung Won Chung official deck, p.~47}

\lecturefigure{hyung-slide-048.jpg}{Comparación tras el paso tres: la final-layer cross-attention pasa a ser una conexión entre capas del mismo nivel}{Hyung Won Chung official deck, p.~48}

\begin{knowledgebox}{Lectura de la figura: por qué el layer index puede representar la granularity}
Las redes profundas suelen formar representaciones jerárquicas: las capas inferiores tienden hacia la forma local y las superiores hacia la semántica abstracta. Si la primera capa del decoder solo puede leer la capa final del encoder, tiene que recuperar información de bajo nivel a partir de representaciones muy abstractas; la conexión entre capas del mismo nivel permite que cada etapa lea entradas de granularidad similar. Esta explicación es una intuición de diseño; si de verdad se forma un cuello de botella todavía requiere una ablation en función de la profundidad.
\end{knowledgebox}

\subsection{Paso cuatro: cambiar la self-attention del encoder a causal}

Los tres primeros pasos ya unificaron el module, los parámetros y la conexión entre capas; la principal diferencia restante es si dentro de la secuencia de entrada hay visibilidad bidireccional. Al cambiar la mask del encoder a causal, los tokens anteriores ya no leen los posteriores; al concatenar la entrada y el objetivo, el grafo de cómputo de las dos stacks se aproxima al de una sola stack decoder-only.

\lecturefigure{hyung-slide-050.jpg}{Paso cuatro: cambiar la bidirectional self-attention del encoder a causal}{Hyung Won Chung official deck, p.~50}

\lecturefigure{hyung-slide-051.jpg}{Tabla de diferencias tras los cuatro pasos: entre ambas arquitecturas solo quedan pequeñas diferencias de implementación}{Hyung Won Chung official deck, p.~51}

\begin{lstlisting}[caption={Pseudocódigo conceptual de la architecture transformation en cuatro pasos}]
model = encoder_decoder()
share(model.cross_attention, model.self_attention)
share(model.encoder_layers, model.decoder_layers)
connect_decoder_layer_to_matching_encoder_layer(model)
set_encoder_mask(model, mode="causal")
concatenate_input_and_target(model)
\end{lstlisting}

\teachervoice{En 00:26:21--00:26:45, Hyung Won sostiene que, tras completar los cuatro pasos, la diferencia entre ambas arquitecturas con la misma tarea y los mismos datos podría ser menor que el ruido de repetir el entrenamiento. Es un juicio empírico, no un teorema formal presentado en el curso; una verificación real todavía exige controlar el número de parámetros, la inicialización, el presupuesto de entrenamiento y el orden de los datos.}

\subsection{Resumen de la sección}

La transformación en cuatro pasos reduce las etiquetas de arquitectura a cuatro hipótesis locales: attention role, separación de parámetros, conexión entre capas y direccionalidad de la entrada. Todas pueden someterse a experimentos por separado y el valor de cada una puede cambiar con la escala. La siguiente sección explicará, una por una, por qué estas estructuras eran razonables en la traducción automática y en las primeras tareas, y por qué deben reexaminarse en los modelos de lenguaje de propósito general, el diálogo de varios turnos y las redes más profundas.
\section{Por qué la estructura funcionó en su momento: traducción e Instruction Tuning como experimento natural}

La sección anterior completó la transformación formal; esta sección pregunta por la motivación histórica. La estructura adicional del encoder-decoder no era una complicación arbitraria: en 2017 la tarea central era la traducción automática, en la que la entrada y el objetivo provienen de idiomas distintos; más tarde, los datos académicos de instruction-tuning presentaban con frecuencia entradas largas y respuestas cortas. Cuando la estructura coincide con la distribución de los datos, el sesgo inductivo eleva de manera notable la eficiencia con un presupuesto limitado.

\subsection{Tres estructuras adicionales}

Esta subsección presenta primero la lista completa y después se centra en el primer elemento. En comparación con decoder-only, el encoder-decoder supone que la entrada y el objetivo son lo bastante distintos, que el objetivo debe leer una entrada ya codificada por completo y que los token de la entrada se prestan a una interacción totalmente bidireccional. Estos tres supuestos corresponden, respectivamente, a los parámetros, a las conexiones entre capas y a la attention mask.

\lecturefigure{hyung-slide-052.jpg}{Las tres estructuras adicionales del encoder-decoder respecto a decoder-only}{Hyung Won Chung official deck, p.~52}

\lecturefigure{hyung-slide-053.jpg}{Primer supuesto: la entrada y el objetivo difieren tanto que justifican parámetros independientes}{Hyung Won Chung official deck, p.~53}

\begin{importantbox}{Tres preguntas para juzgar si un inductive bias sigue siendo válido}
\begin{enumerate}
  \item ¿Con qué distribución de datos o restricción de la tarea coincidía originalmente?
  \item ¿Los datos, el objetivo y la forma de interacción actuales siguen cumpliendo esa restricción?
  \item Si se retira la estructura y se aumenta la escala, ¿en qué punto del presupuesto se cruzan las curvas de desempeño?
\end{enumerate}
Responder solo la primera pregunta lleva a confundir la validez histórica con una corrección permanente.
\end{importantbox}

\subsection{Traducción automática: los parámetros independientes eran naturales}

En 2017, el Transformer se orientaba a la source-to-target translation. La entrada es una oración en el idioma fuente y el objetivo es la salida en otro idioma; ambos difieren en vocabulario, orden de palabras y papel en la generación. En un escenario que solo optimiza la traducción, hacer que el encoder se especialice en el idioma fuente y el decoder en el idioma objetivo es un supuesto previo eficiente. Esta subsección reconoce primero esa validez histórica y después compara qué condiciones cambió el preentrenamiento general, para no deducir desde las interfaces de tareas actuales que el diseño anterior era «erróneo desde el principio».

\lecturefigure{hyung-slide-054.jpg}{En la traducción automática, la diferencia de distribución entre source y target respalda los parámetros independientes}{Hyung Won Chung official deck, p.~54}

\lecturefigure{hyung-slide-055.jpg}{Un modelo de lenguaje general debe integrar conocimiento entre idiomas, lo que debilita la motivación de los parámetros independientes}{Hyung Won Chung official deck, p.~55}

Al leer la segunda diapositiva conviene notar que la tarea ya cambió. El preentrenamiento moderno no solo aprende traducción del inglés al alemán, sino que debe formar un conocimiento del mundo compartido a partir de texto en muchos idiomas; un mismo hecho puede expresarse en idiomas distintos. Si la entrada y el objetivo quedan aislados de forma permanente, al modelo puede costarle más reutilizar sus representaciones. Por otro lado, la especialización multilingüe todavía puede tener valor, de modo que esta diapositiva plantea un scale-era counterargument, no que los parámetros independientes sean del todo inútiles.

\teachervoice{La comparación histórica de 00:27:27--00:28:45 es muy concreta: cuando el objetivo es solo traducir, separar los parámetros por idioma es «muy natural»; cuando el objetivo pasa a ser el aprendizaje de conocimiento general, el ponente prefiere fusionar el conocimiento entre idiomas. La evaluación de un inductive bias cambia junto con el objective.}

\subsection{FLAN: entradas largas y objetivos cortos dan una ventaja inesperada al encoder-decoder}

Instruction tuning ofrece un experimento natural más fino. El mismo equipo aplicó fine-tuning a PaLM (decoder-only) y a T5 (encoder-decoder); aunque dedicó la mayor parte del tiempo a optimizar PaLM, la ganancia de T5 fue mayor. El ponente no atribuyó el resultado a una superioridad general de la arquitectura, sino que volvió a la distribución de longitudes de los datos para buscar la correspondencia.

\lecturefigure{hyung-slide-056.jpg}{Instruction finetuning unifica una gran cantidad de tareas académicas como instrucciones en lenguaje natural}{Hyung Won Chung official deck, p.~56}

\lecturefigure{hyung-slide-057.jpg}{En el experimento FLAN, la ganancia del fine-tuning es mayor para el encoder-decoder}{Hyung Won Chung official deck, p.~57}

\begin{knowledgebox}{Lectura de la figura: se compara la ganancia, no qué modelo gana en términos absolutos}
La tabla destaca la improvement antes y después del instruction tuning. No basta con mirar el recuadro rojo para concluir que T5 supera a PaLM en todas las tareas y en todas las escalas; la conclusión de la clase es que la estructura del encoder-decoder coincide especialmente con ese conjunto de datos académicos, por lo que aprovecha mejor la señal del fine-tuning.
\end{knowledgebox}

\lecturefigure{hyung-slide-058.jpg}{Los datos académicos suelen tener una distribución de longitudes con entradas largas y objetivos cortos}{Hyung Won Chung official deck, p.~58}

Si la longitud de la entrada es la variable aleatoria $L_x$ y la longitud del objetivo es $L_y$, este conjunto de datos cumple aproximadamente
\[
\mathbb{E}[L_x]\gg \mathbb{E}[L_y].
\]
Aquí, $\mathbb{E}[L_x]$ y $\mathbb{E}[L_y]$ son las longitudes promedio de la entrada y del objetivo, respectivamente. Los enunciados largos y las etiquetas cortas hacen que las distribuciones de ambos lados difieran claramente, y los parámetros independientes de encoder/decoder pueden especializarse; pero la generación de textos largos y el chat de varios turnos reducen e incluso invierten esa asimetría.

\teachervoice{En 00:29:22--00:30:46, Hyung Won dice que el equipo dedicó el 99\% del tiempo a optimizar PaLM y, sin embargo, T5 obtuvo una ganancia mayor en pocos días; su explicación es que el sesgo de longitud de los datos y la arquitectura «coincidieron por casualidad». Este teacher voice convierte el resultado de una comparación de marcas en un análisis de correspondencia de distribuciones.}

\begin{warningbox}{La facilidad de evaluar un academic benchmark moldea la distribución de entrenamiento}
Los objetivos largos son difíciles de anotar manualmente y de calificar en forma automática, por lo que los datos académicos se inclinan hacia entradas largas y respuestas cortas. Que algo sea fácil de calificar no equivale a que tenga valor para el usuario: la escritura creativa, el código, el análisis y la conversación suelen requerir salidas largas. Si la arquitectura del modelo se optimiza solo para la distribución del benchmark, la ventaja estructural original puede desaparecer cuando cambian las tareas de despliegue.
\end{warningbox}

\subsection{Resumen de la sección}

La traducción automática y FLAN muestran por qué el sesgo inductivo puede producir beneficios reales en una época determinada: dirige una capacidad limitada hacia lo que la distribución de los datos más necesita. Pero el preentrenamiento general, la generación larga y el diálogo de varios turnos cambiaron la relación entre entrada y objetivo, y obligaron a los investigadores a volver a probar la separación de parámetros. La historia no sirve para ridiculizar los diseños anteriores, sino para identificar las condiciones en las que un diseño es válido.
\section{Granularidad de la representación, bidirectionality y caché en múltiples turnos}

La sección anterior analizó la separación de parámetros; esta sección aborda las dos estructuras restantes: que el decoder solo lea la capa final del encoder y que el encoder use bidirectional attention. Ambas mejoraron en su momento la calidad en las tareas, pero a medida que las redes se vuelven más profundas y las conversaciones más largas, los problemas de granularidad de la información y de cómputo repetido se vuelven notorios. Aquí se colocan por primera vez la architecture choice y el serving cost en una misma figura.

\subsection{¿Leer solo la última encoder layer crea un cuello de botella de información?}

Esta subsección se centra en el punto desde el que lee la cross-attention. En las redes profundas, las distintas capas suelen representar información de distinta granularidad; si incluso la primera capa del decoder solo puede acceder a la capa final del encoder, la información formal de las capas bajas puede quedar excesivamente comprimida. En los modelos actuales de unas decenas de capas el efecto quizá sea pequeño, pero una red extremadamente profunda podría amplificar este desajuste.

\lecturefigure{hyung-slide-059.jpg}{Segunda suposición: el objetivo solo necesita leer una entrada ya completamente codificada}{Hyung Won Chung official deck, p.~59}

\lecturefigure{hyung-slide-060.jpg}{La cross-attention estándar solo lee las activaciones de la capa final del encoder}{Hyung Won Chung official deck, p.~60}

\lecturefigure{hyung-slide-061.jpg}{En las redes profundas, las capas inferiores y superiores suelen codificar información de distinta granularidad}{Hyung Won Chung official deck, p.~61}

\begin{knowledgebox}{Lectura de la figura: la representation hierarchy es una heurística, no una tabla semántica fija}
En las redes de visión se suele ilustrar la representación jerárquica con «bordes en las capas bajas, objetos en las capas altas»; en los modelos de lenguaje también puede aparecer una tendencia que va de la forma local a la semántica abstracta. Sin embargo, el significado de cada capa cambia con el objetivo de entrenamiento, las conexiones residuales y la profundidad. Lo que la figura realmente sustenta es que distintas capas pueden aportar información complementaria, por lo que exponer solo la capa final es una decisión estructural que debe someterse a prueba.
\end{knowledgebox}

\teachervoice{00:32:19--00:33:01, Hyung Won dice que, en los T5 de alrededor de 24 capas que él ha usado, que la layer 1 lea la layer 24 parece no causar problemas; pero si en el futuro hubiera diez o mil veces más profundidad, él «no se sentiría del todo cómodo». Es una pregunta típica de scale extrapolation, no una falla ya observada.}

\subsection{Qué contrafáctico representa la layer-aligned connection}

La subsección anterior señaló el problema de granularidad; esta construye la alternativa mínima: que la capa $\ell$ del decoder lea la capa $\ell$ del encoder. Así, cada etapa recibe una representación de entrada de profundidad similar, lo que también se acerca más a la token interaction en la misma capa de un decoder-only. El valor del contrafáctico está en que permite experimentos controlados, aunque al final no se adopte.

\lecturefigure{hyung-slide-062.jpg}{Tercera suposición: dentro de la entrada se necesita all-to-all bidirectional interaction}{Hyung Won Chung official deck, p.~62}

\lecturefigure{hyung-slide-063.jpg}{Arquitecturas alternativas con layer-aligned connection y causal input}{Hyung Won Chung official deck, p.~63}

\begin{importantbox}{El producto correcto del análisis estructural es un ablation axis}
De H059--H063 se obtienen dos ejes experimentales independientes: final-layer versus layer-aligned cross-attention, y bidirectional versus causal input mask. Si se cambian ambos a la vez, la diferencia de desempeño no puede atribuirse a ninguno; la derivación en cuatro pasos de la clase sirve precisamente para generar ablations que puedan escalarse por separado.
\end{importantbox}

\subsection{Beneficio en calidad y costo de sistema de la bidirectionality}

En la época de BERT, el bidirectional encoder era muy importante para las tareas difíciles de comprensión, porque cada token puede aprovechar al mismo tiempo el contexto izquierdo y el derecho. Sin embargo, en el instruction tuning a gran escala, la anecdotal experience del ponente es que la diferencia entre bidireccional y unidireccional es pequeña; al mismo tiempo, el chat de múltiples turnos expone un costo de sistema evidente: cada mensaje nuevo cambia el «futuro» de la entrada anterior, y la representación bidireccional debe volver a codificarse.

\lecturefigure{hyung-slide-064.jpg}{A escala suficiente, el beneficio de la bidirectionality puede reducirse, pero aumenta el costo de la conversación de múltiples turnos}{Hyung Won Chung official deck, p.~64}

\teachervoice{00:33:11--00:33:59, Hyung Won marca explícitamente la conclusión como highly anecdotal: en Flan-2, la diferencia entre el fine-tuning bidireccional y el unidireccional es pequeña. Estas notas conservan ese nivel de evidencia; no debe reescribirse como «se ha demostrado que la bidirectional attention es inútil».}

\subsection{Por qué la causal attention puede reutilizar la KV cache}

Esta subsección convierte la figura en un balance de cómputo. En un causal model, la representación de los tokens anteriores no depende de los tokens nuevos futuros, por lo que sus key/value calculados en el turno previo pueden guardarse en caché; al añadir un mensaje solo se calcula la parte nueva. En un bidirectional encoder, los tokens anteriores pueden leer los tokens nuevos, de modo que tras añadirlos las representaciones previas dejan de ser válidas y normalmente hay que volver a codificar toda la entrada.

\lecturefigure{hyung-slide-066.jpg}{Comparación entre la recodificación bidirectional y la cache reuse unidirectional en una conversación de múltiples turnos}{Hyung Won Chung official deck, p.~66}

Sea $n$ la longitud del contexto existente y $m$ la longitud añadida. Ignorando constantes y el MLP, el costo de rehacer la full attention es aproximadamente
\[
O\bigl((n+m)^2\bigr),
\]
mientras que la causal KV cache (key-value cache, que guarda los tensores key/value de cada capa para los tokens anteriores) solo necesita calcular, para las nuevas query, la attention sobre la caché anterior y los tokens añadidos, con un costo incremental aproximado de
\[
O(nm+m^2).
\]
Aquí $n$ es la longitud del contexto anterior y $m$ es el número de tokens añadidos. La caché reduce el cómputo repetido, pero aumenta el uso de memoria de la GPU; el servicio con contextos largos todavía requiere gestionar cache placement, batching y eviction.

\begin{knowledgebox}{Lectura de la figura: la región azul reutilizable corresponde a la invariancia causal}
En el caso unidirectional, los tokens anteriores no pueden ver el futuro, por lo que la llegada de un mensaje nuevo no cambia sus key/value; en el caso bidirectional, la representación de ``How are you?'' cambia porque después aparece ``Bad''. La ventaja de eficiencia del sistema proviene de la dirección de las dependencias en el grafo de cómputo, no del nombre decoder-only en sí.
\end{knowledgebox}

\teachervoice{La explicación en clase de 00:34:07--00:35:12 conecta la calidad del modelo con el serving: una estructura que en 2018 mejoró los benchmark puede generar costos de recodificación en la forma de producto de múltiples turnos. La architecture trajectory no la determina solo la accuracy, sino también el patrón de interacción.}


\subsection{Conclusión: convertir el análisis histórico en la siguiente ablación}

La última diapositiva no pide a los estudiantes copiar una arquitectura fija, sino volver a su propio problema de investigación: ¿qué estructuras predeterminadas hay en el sistema actual? ¿Por qué eran útiles en el régimen anterior? Al cambiar el cómputo, los datos, la profundidad de la red y la longitud de la interacción, ¿pasarán de ser un atajo a ser un cuello de botella?

\lecturefigure{hyung-slide-067.jpg}{Conclusión: identificar la fuerza dominante y analizar la estructura adicional desde la perspectiva del scaling}{Hyung Won Chung official deck, p.~67}

\teachervoice{00:35:42--00:36:18, al cerrar, Hyung Won pide a los estudiantes que hagan análisis histórico una y otra vez, para entrenarse en identificar «qué suposiciones hay que revisar, por qué, y si pueden sustituirse por un mecanismo más general y hacer scale up». Esto es más importante que recordar que ganó decoder-only.}

\begin{importantbox}{De la conclusión histórica a la pregunta experimental}
Para cualquier estructura $b$, escriba al menos cuatro elementos: el problema anterior que resolvía, su beneficio actual, el eje de escala que podría limitar y la ablation capaz de aislar esa suposición. Si solo se sabe que «al quitarla baja la puntuación actual», todavía no se ha respondido si cambiará la scaling trajectory a largo plazo.
\end{importantbox}

\subsection{Resumen de la sección}

La jerarquía de representaciones nos recuerda revisar el final-layer bottleneck; la interacción de múltiples turnos nos recuerda considerar la attention direction junto con el costo de la caché; el cierre condensa estos casos en un procedimiento general de investigación: identificar las suposiciones predeterminadas, explicar su beneficio en el entorno anterior y, a lo largo del nuevo eje de escala, comprobar si todavía vale la pena conservarlas.
\section{Síntesis y ampliación}

Hyung Won Chung (钟亨原) completó con un deck de 67 páginas una «arqueología de arquitecturas»: primero planteó la fuerza direccional de un exponentially cheaper compute; después usó la Bitter Lesson para explicar la tensión entre estructura y escala; luego descompuso el paso de encoder-decoder a decoder-only en cuatro transformaciones locales, y por último recurrió a la traducción, a FLAN, a las representaciones jerárquicas y a la caché en conversaciones de varios turnos para explicar por qué cada sesgo fue razonable en su momento y por qué puede haber caducado.

\subsection{Diez conclusiones centrales}

\begin{enumerate}
  \item Al estudiar un campo que cambia con rapidez, además de seguir los resultados hay que estudiar la trayectoria del cambio.
  \item La dominant driving force es una aproximación de modelado con márgenes de error, no una declaración de causa única.
  \item La tendencia histórica de compute per dollar puede orientar la estrategia de investigación, pero no demuestra por sí sola la pendiente futura.
  \item El inductive bias aporta a la vez eficiencia presente y un posible scale ceiling.
  \item ``Scalable'' significa que el desempeño sigue mejorando a medida que crece el presupuesto; no equivale a poder ejecutarse de forma distribuida.
  \item Encoder-decoder, encoder-only y decoder-only comparten la base de sequence-model, y sus diferencias pueden descomponerse localmente.
  \item Las cuatro transformation retiran, respectivamente, los supuestos estructurales sobre los parámetros de attention, los parámetros del stack, la layer connection y la mask.
  \item La traducción automática y los datos de long-input/short-target explican por qué encoder-decoder tuvo ventaja en su momento.
  \item El valor tanto de la final-layer-only cross-attention como de la bidirectionality debe volver a medirse según la profundidad y el modo de interacción.
  \item Lo que realmente ofrece la historia de las arquitecturas es un bias audit, no una lista eterna de ganadores.
\end{enumerate}

\subsection{Una tabla que condensa el ciclo de vida de las estructuras}

\begin{table}[H]
\centering
\small
\begin{tabular}{@{}L{0.17\textwidth}L{0.24\textwidth}L{0.24\textwidth}L{0.25\textwidth}@{}}
\toprule
Supuesto estructural & Justificación en el entorno anterior & Cambio de regime & Prueba que conviene hacer \\
\midrule
Separar los parámetros de entrada y de objetivo & En la traducción, los dos lados tienen papeles y vocabularios distintos & Conocimiento general, generación larga, diálogo de varios turnos & Scaling curve con parámetros compartidos \\
Leer solo la última encoder layer & Suficiente con profundidad moderada y sencillo de implementar & Redes más profundas, jerarquías de granularidad más ricas & Layer-aligned cross-attention \\
Attention bidireccional sobre la entrada & Tareas de comprensión sobre la entrada completa & Streaming e interacción incremental de varios turnos & Ganancia de calidad frente al costo de recodificar \\
Architecture label fija & Facilita la implementación y la comparación & Componentes sustituibles localmente, cambios en las restricciones del sistema & Ablation de una sola variable por transformation \\
\bottomrule
\end{tabular}
\caption{Expresar las estructuras históricas como hipótesis de ciclo de vida que pueden verificarse.}
\end{table}

\subsection{Práctica: el procedimiento mínimo de un bias audit}

El siguiente pseudocódigo convierte el método de la clase en un registro de investigación. No sustituye a los experimentos, pero evita que el equipo reporte solo la mejor puntuación actual y olvide registrar las condiciones en las que la estructura es válida.

\begin{lstlisting}[caption={Plantilla de registro de investigación para un inductive-bias audit}]
for bias in architecture.assumptions:
    document(original_problem(bias))
    measure(current_regime_gain(bias))
    choose_scale_axes(compute, data, depth, interaction)
    run_single_variable_ablation(bias, scale_axes)
    report(crossover_point, uncertainty, system_cost)
\end{lstlisting}

\begin{lstlisting}[caption={Pseudocódigo de decisión para distinguir el óptimo actual de la tendencia de largo plazo}]
if deployment_budget < observed_crossover:
    ship(currently_efficient_method)
else:
    validate(less_structured_method)
keep_both_curves_and_failure_cases()
never_label_extrapolation_as_measurement()
\end{lstlisting}

\begin{warningbox}{No convertir la «investigación por sustracción» en un nuevo dogma}
Retirar estructura no es, por naturaleza, más científico. Si un sesgo expresa una invariancia real, reduce de forma notable la necesidad de datos o cumple restricciones de latencia, privacidad y consumo de energía, puede seguir siendo razonable a largo plazo. Lo correcto es medir beneficios, costos y el punto de cruce, no tratar ``less structure'' como una preferencia estética que no requiere evidencia.
\end{warningbox}

\subsection{Preguntas de autoevaluación}

\begin{enumerate}
  \item ¿Por qué el experimento de dejar caer una pluma permite ignorar la resistencia del aire, pero no permite concluir que un sistema complejo tiene siempre una sola variable dominante?
  \item ¿En qué se diferencian compute per dollar y la Moore's law? ¿Qué es lo que la figura de la clase no demuestra?
  \item ¿Cómo puede el inductive bias aumentar la sample efficiency y, al mismo tiempo, reducir la escalabilidad a largo plazo?
  \item ¿Cuáles son las tres attention role del encoder-decoder?
  \item ¿Qué supuesto estructural retira cada uno de los cuatro pasos de transformation?
  \item ¿Por qué la traducción automática favorece de forma natural la separación de parámetros de entrada y de objetivo?
  \item ¿Por qué la mayor ganancia de FLAN-T5 solo respalda la hipótesis sobre la distribución de los datos, en lugar de constituir una prueba causal?
  \item ¿Qué cuello de botella puede provocar la final-layer-only cross-attention en una red muy profunda?
  \item ¿Por qué la bidirectional attention impide reutilizar la KV cache en un diálogo de varios turnos?
  \item Si la tendencia del compute cambia, ¿qué supuestos deben actualizarse, en lugar de conservar mecánicamente cuál conclusión?
\end{enumerate}

\subsection{Lecturas adicionales}

\begin{itemize}
  \item Richard Sutton, \emph{The Bitter Lesson}: la relación de largo plazo entre los métodos generales, el cómputo y la estructura diseñada por personas.
  \item Vaswani et al., \emph{Attention Is All You Need}: el Transformer encoder-decoder original y los tres tipos de attention role.
  \item Devlin et al., \emph{BERT}: el diseño representativo de encoder-only y del masked language modeling.
  \item Brown et al., \emph{Language Models are Few-Shot Learners}: un hito histórico del scaling de decoder-only.
  \item Chung et al., \emph{Scaling Instruction-Finetuned Language Models}: la evidencia de instruction-tuning de FLAN-T5/PaLM.
\end{itemize}

\end{document}
